%-----------------------------------------------------------
% Wire Guides
% \label{subsec:wire_guides}
%-----------------------------------------------------------

A network of internal wire guides was designed to route and align the antenna conductors within the gel volume. 
Each guide follows a unique, non-overlapping path leading to a carefully selected endpoint inside the cavity. 
The endpoints were positioned to intentionally break possible spatial symmetries in electrical stimulation, 
thereby enabling more representative and diverse field distributions during testing.

All guides maintain a uniform circular cross-section of \nota{DIAMETER HERE} 
to accommodate a range of cable or electrode diameters. 
This geometry also facilitates easy insertion and sealing of the antennas by any preferred method, 
without requiring a predefined connector or cap design. 
The open-ended approach ensures compatibility with various electrode terminals and sealing techniques.

Each guide is uniquely shaped and numbered at its base for unambiguous identification 
during assembly. 
The numbering system allows consistent correspondence with the experimental layout, \nota{have to decide if we put the numbers on the lid, or provide a diagram for this -- but the guides are id-ed}

The overall design of the wire-guide network enables:
\begin{itemize}
    \item straightforward insertion of flexible antenna wires,
    \item precise alignment of electrode tips within the conductive medium,
    \item and reliable mechanical fixation during gel solidification.
\end{itemize}

\begin{figure}[h]
    \centering
    \includegraphics[width=0.65\textwidth]{figures/wire_guides.png}
    \caption{Wire guides integrated within the head shell, showing unique routing and numbered bases.}
    \label{fig:wire_guides}
\end{figure}

\nota{Agregar fotos del ensemble + cada una + numero + connector maybe}
\nota{Add table listing each wire guide ID, type, and target location.}
